\documentclass[11pt]{article}
\usepackage[margin=1in]{geometry}
\usepackage{amsmath,amssymb,amsthm,mathtools}
\usepackage[hidelinks,pdfusetitle]{hyperref}
\usepackage{microtype}
\usepackage{enumitem}
\usepackage{array}
\newcommand{\DOI}[1]{\mbox{\href{https://doi.org/#1}{\nolinkurl{#1}}}}
\newtheorem{theorem}{Theorem}[section]
\newtheorem{lemma}[theorem]{Lemma}
\newtheorem{corollary}[theorem]{Corollary}
\newtheorem{proposition}[theorem]{Proposition}
\theoremstyle{definition}
\newtheorem{definition}[theorem]{Definition}
\theoremstyle{remark}
\newtheorem{remark}[theorem]{Remark}
\DeclareMathOperator{\Tr}{Tr}
\DeclareMathOperator{\spec}{spec}
\DeclareMathOperator{\dist}{dist}
\DeclareMathOperator{\diag}{diag}


\newcommand{\E}{\mathbb{E}}
\newcommand{\R}{\mathbb{R}}
\newcommand{\C}{\mathbb{C}}
\newcommand{\eff}{\mathrm{eff}}
\newcommand{\pair}{\mathrm{pair}}
\newcommand{\Ggreen}{\mathcal{G}}
\newcommand{\Tpair}{T_{\pair}}
\newcommand{\Tlin}{T_{\mathrm{lin}}}
\newcommand{\Gb}{G_{\mathrm{b}}}
\newcommand{\Ueff}{U_{\eff}}
\emergencystretch=1.5em
\title{Method obstructions to linear-in-frequency pairing ceilings in
electron-phonon models: integrated cumulants, positive determinants, and the
vanishing effective attraction}
\author{Davis Rens\\SuperMatics Inc.\\\href{https://orcid.org/0000-0001-5839-5197}{ORCID 0000-0001-5839-5197}}
\date{}
\newcommand{\PAPERDOI}{10.5281/zenodo.24408871}
\newcommand{\SUPPDOI}{10.5281/zenodo.24408875}
\newcommand{\REPOURL}{https://github.com/SuperMatics-Inc/linear-pairing-ceiling-obstructions}
\begin{document}
\maketitle

\begin{center}
\small
Article DOI: \DOI{\PAPERDOI}\\
Supplement DOI: \DOI{\SUPPDOI}\\
Source repository: \url{\REPOURL}
\end{center}

\begin{abstract}
Can the pair-ordering temperature of a Holstein or bond-SSH electron-phonon model
be bounded at a temperature scale proportional to the phonon frequency, uniformly
in that frequency and at fixed hopping and coupling? That linear step remains open,
and this note does not close it. It records three exact obstructions that identify
which proof methods cannot reach it, together with one case where a
frequency-uniform ceiling does survive. First, for the Holstein model with
arbitrary real static site fields, the exact first coefficient of the local
dynamic cumulant, computed with the dynamic Gaussian already integrated out,
is at least $U\Omega^2/(1200\tau^3)$ on an open three-site path for
$0<\Omega\le\tau\le t/4$, and at least $U\Omega^2/(2400\tau^3)$ on every square
torus of side at least four with a finite static barrier. At $\tau=C\Omega$ this
diverges as $\Omega\downarrow0$, so no expansion whose uniform local activity must
dominate it can stay small. Second, a strictly positive single-impurity
determinant does not bound the impurity propagator: on a two-slice history the
Fermion determinant exceeds one, all monodromy eigenvalues are positive, and the
conditional Green function grows as $B^2/(2v^2)$. An interval-certified fixture at
$B=6$, $v=1/2$ gives $\det(1+U)>15$ and $|\Ggreen(1,0)_{00}|>52.89$. Third, for the
pure bond-SSH model with vanishing effective attraction the high-temperature
regime is closed: for $\beta|t|+\beta\Gb\le2^{-10}$ the pair correlator is at most
$(341/103)(32/341)^{d}$, uniformly in frequency, chemical potential and density, so
$k_B\Tpair\le1024(|t|+\Gb)$, and three further proof routes are obstructed by exact
identities. Nothing here is a bound on a measured transition temperature, and no
materials claim is made.
\end{abstract}

\section{Question, scope, and prior work}

We study the spinful Holstein and bond-SSH models on square tori, at fixed
nearest-neighbour hopping $t>0$, fixed electron-phonon coupling, fixed density or
chemical potential, and phonon frequency $\Omega>0$. In the Holstein case each site
carries an independent harmonic oscillator whose displacement couples linearly to
the local density, with relaxation energy $U>0$. In the bond-SSH case each
nearest-neighbour bond carries an independent oscillator whose displacement
modulates the bond hopping amplitude, with bond relaxation energy $\Gb=g^2/K$ and
vanishing local attraction. Units are $k_B=\hbar=1$, with $\tau=\beta^{-1}=k_BT$
written as an energy; physical temperatures are restored only when a model scale is
compared with $k_BT$.

The question that this series of notes addresses is the following. At fixed $t$ and
fixed coupling, and at regular density, can one bound pair correlations at a
temperature scale proportional to the phonon frequency, uniformly in
$\Omega\downarrow0$? A positive answer would be a pair ceiling of the form
$k_B\Tpair\le C(t,U)\,\Omega$ with $C$ independent of $\Omega$. Proof proposals for
that linear scale are numerous, and several weaker ceilings are established:
frequency-uniform bounds at a constant temperature scale in the strong-coupling
regime \cite{stage6}, a two-thirds-power adiabatic bound and its sharpness
\cite{linear4,linear5}, slow-phonon estimates \cite{slow}, an explicit pair bound
for the attractive Hubbard model with bond SSH coupling \cite{sshbond3}, method
obstructions from single-spin determinant zeros \cite{barriers}, and limits of
defect-insertion arguments \cite{limits}.

This note is the second collection of method obstructions in that programme. The
first \cite{barriers} studied determinant zeros of single-spin propagators and the
limits of specified cutoff and twist families. The present note adds three further
obstructions of a different character, and one positive result.

\begin{enumerate}[label=(\Alph*),leftmargin=2.2em]
\item The dynamic Gaussian can be integrated out first, around the strictly
positive static free-fermion trace, with no static conditioning and no determinant
zero. Even then, the exact first local dynamic cumulant is not uniformly small at a
linear temperature scale: an explicit three-site static field at fixed bare hopping
produces a resonance whose contribution is $\Theta(U\Omega^2/\tau^3)$. This blocks
unresummed expansions about the static reference whose local activity must dominate
that cumulant (Section~\ref{sec:cumulant}).
\item A positive Fermion determinant, even together with positive monodromy
eigenvalues, does not bound a conditional impurity Green function. An exact
two-slice family gives $\det(1+U)>1$ and $|\Ggreen(1,0)_{00}|\sim B^2/(2v^2)$, so no
pathwise history-independent propagator bound follows from determinant positivity
alone (Section~\ref{sec:det}).
\item For the pure bond-SSH model at $\Ueff=0$, the high-temperature regime is
closed by a frequency-uniform polymer expansion, and three further routes are
obstructed by exact identities: no phonon-induced instantaneous atomic gap, a
non-terminating bond Lang-Firsov expansion, and free-fermion algebraic decay that
forbids any hopping-uniform exponential clustering (Section~\ref{sec:survives}).
\end{enumerate}

None of these statements is a physical counterexample. Each is a statement about a
precisely described class of arguments applied to a precisely described model
quantity. In particular, none of them shows that $\Tpair/\Omega$ can diverge, and
none shows that a linear ceiling is false. The linear step remains open, and
Section~\ref{sec:open} lists what is missing.

Section~\ref{sec:model} fixes the models and defines the model quantity
$\Tpair$ that every result in this note concerns. Sections~\ref{sec:cumulant},
\ref{sec:det} and \ref{sec:survives} contain the results.
Section~\ref{sec:notclaimed} states what is not claimed, and
Section~\ref{sec:open} lists open problems.

\section{Models, notation, and the shape of a linear ceiling}
\label{sec:model}

\subsection{The Holstein model}

Let $\Lambda_L=(\mathbb{Z}/L\mathbb{Z})^2$ be a square torus of side $L\ge4$ with
nearest-neighbour edge set $\mathcal{E}_\Lambda$ and graph distance $d$. Fermion
operators $c_{x\sigma}$ carry spin $\sigma\in\{\uparrow,\downarrow\}$, and
$n_{x\sigma}=c_{x\sigma}^\dagger c_{x\sigma}$, $n_x=n_{x\uparrow}+n_{x\downarrow}$,
$\Delta_x=c_{x\downarrow}c_{x\uparrow}$. Each site carries a coordinate $X_x$ and
momentum $P_x$ with $[X_x,P_{x'}]=i\delta_{xx'}$, mass $M_x$, spring constant
$K_x>0$ and frequency $\Omega_x=\sqrt{K_x/M_x}$. With a linear force constant
$F_x>0$, the Hamiltonian is
\begin{equation}
H_\Lambda=H_{\mathrm{el}}+\sum_{x\in\Lambda_L}\left(\frac{P_x^2}{2M_x}
+\frac{K_x}{2}X_x^2+F_xX_x\,(n_x-1)\right)-t\sum_{\langle xy\rangle,\sigma}
\left(c_{x\sigma}^\dagger c_{y\sigma}+c_{y\sigma}^\dagger c_{x\sigma}\right),
\label{eq:H}
\end{equation}
where the electronic part may also contain an arbitrary real static site field
$W_x$ (including the chemical potential). Two normalisations are useful and we use
both:
\begin{equation}
U_x:=\frac{F_x^2}{K_x},\qquad
F_xX_x=G_x\,b_x,\quad b_x=a_x+a_x^\dagger,\quad G_x=\sqrt{\frac{U_x\Omega_x}{2}}.
\label{eq:U}
\end{equation}
Thus $U_x$ is the site relaxation energy and $G_x$ is the boson-operator coupling of
the site, with $G_x^2=U_x\Omega_x/2$. The double commutator of the force-scaled
displacement with the Hamiltonian is
\begin{equation}
\bigl[F_xX_x,[H_\Lambda,F_xX_x]\bigr]=\frac{F_x^2}{M_x}=U_x\Omega_x^2 .
\label{eq:dcomm}
\end{equation}
Where a single coupling is quoted we write $U=F^2/K$ and $G^2=U\Omega/2$ for a
uniform model. Earlier notes
in this series sometimes write the site coupling as $U=g^2/K$; that lowercase $g$ is
the force constant $F$ of \eqref{eq:H}, while here lowercase $g$ is reserved for the
bond force constant of the Peierls model \eqref{eq:Hssh}. The mass $M_x$ and the
spring constant $K_x$ enter the proofs only through $\Omega_x$ and $U_x$.

\subsection{The pure bond-SSH model}

For the bond-SSH model each bond $b=\{x,y\}\in\mathcal{E}_\Lambda$ carries
$X_b,P_b$ with mass $M$, spring constant $K>0$ and $\Omega=\sqrt{K/M}$, and
\begin{equation}
H_\Lambda=-\mu N-t\sum_{b\in\mathcal{E}_\Lambda}T_b
+g\sum_{b\in\mathcal{E}_\Lambda}X_bT_b+\sum_{b\in\mathcal{E}_\Lambda}
\left(\frac{P_b^2}{2M}+\frac{K}{2}X_b^2\right),\qquad
T_b=\sum_\sigma\left(c_{x\sigma}^\dagger c_{y\sigma}+c_{y\sigma}^\dagger
c_{x\sigma}\right),
\label{eq:Hssh}
\end{equation}
so that the local effective attraction is $\Ueff=0$ and the only coupling is the
Peierls one. We write
\begin{equation}
\Gb:=\frac{g^2}{K}
\label{eq:Gb}
\end{equation}
for the bond relaxation energy. Here $g$ is a bond force constant and must not be
confused with the site value $G_x$ of \eqref{eq:U}; when both models appear in one
sentence we write the Holstein site coupling out in full. The proved pair bound of
\cite{sshbond3} for the attractive Hubbard model with bond SSH coupling assumes
$\beta t^2/\Ueff\le2^{-27}$ together with $\beta\Gb\le2^{-13}$, with prefactor
$1965840/475579$ and decay base $48645/524224$. At $\Ueff=0$ and $t>0$ that first
hypothesis cannot hold, which is the case that Section~\ref{sec:survives} closes.

\subsection{The model quantity \texorpdfstring{$\Tpair$}{T\_pair}}

Every result in this note concerns the following model quantity and never a measured
transition temperature.

\begin{definition}[pair-decay temperature]
\label{def:tpair}
Fix $t$, the coupling parameters and $\Omega$. Let $\mathcal{Q}$ be the set of
temperatures $T>0$ for which there exist constants $A=A(t,\text{coupling},T)<\infty$
and $q=q(t,\text{coupling},T)<1$, independent of $L\ge4$ and of the chemical
potential, such that
\begin{equation}
\left|\langle\Delta_x^\dagger\Delta_y\rangle_{\Lambda_L}\right|\le A\,q^{d(x,y)}
\label{eq:clustering}
\end{equation}
holds for every $L\ge4$, every $\mu$, and all distinct $x,y\in\Lambda_L$, in the
grand-canonical state at inverse temperature $\beta=1/T$. Set
$\Tpair:=\inf\mathcal{Q}$.
\end{definition}

\begin{remark}
A linear ceiling is the assertion $k_B\Tpair\le C(t,U)\,\Omega$ with $C$ independent
of $\Omega$, that is, exponential decay of \eqref{eq:clustering} for every temperature
above a scale proportional to the frequency. Write $\Tlin=C(t,U)\Omega$. An
$\Omega$-uniform ceiling of the form $k_B\Tpair\le C(t,U)$ with $C$ not vanishing as
$\Omega\downarrow0$ is a different statement: it says nothing at the scale $\Tlin$,
and in particular does not imply a linear ceiling. The results of
Section~\ref{sec:survives} are of the second type.
\end{remark}

\subsection{The route classes examined here}

Four families of argument have been proposed for the linear scale. Their current
status in this programme is summarised in Table~\ref{tab:routes}; the new content of
this note is the entries marked with a reference to the sections below.

\begin{table}[htbp]
\centering
\small
\begin{tabular}{@{}>{\raggedright\arraybackslash}p{0.24\linewidth}>{\raggedright\arraybackslash}p{0.29\linewidth}>{\raggedright\arraybackslash}p{0.37\linewidth}@{}}
\hline
Route & Obstruction & Established here \\
\hline
Condition on static fields, then expand the integrated dynamic Gaussian &
Exact first cumulant is $\Theta(U\Omega^2/\tau^3)$ at fixed hopping &
Theorems~\ref{thm:witness} and~\ref{thm:torus}, Section~\ref{sec:cumulant} \\
\hline
Per-history determinant positivity or nonvanishing, giving a uniform propagator
bound & Positive determinant with arbitrary conditional amplification &
Theorem~\ref{thm:amp}, Corollary~\ref{cor:fixture}, Section~\ref{sec:det} \\
\hline
Ward identities and Schur leakages on the Matsubara grid & Threshold algebra at
$\alpha=1$, diagonal versus operator norm & Remarks~\ref{rem:ward} and
\ref{rem:dimer}, Section~\ref{sec:det} \\
\hline
High-temperature polymer or cluster expansion at $\Ueff=0$ & Not an obstruction;
the route closes & Theorems~\ref{thm:gc} and~\ref{thm:canon},
Section~\ref{sec:survives} \\
\hline
Phonon-induced gap, bond Lang-Firsov, hopping-uniform clustering at $\Ueff=0$ &
Exact identities & Theorem~\ref{thm:routes}, Section~\ref{sec:survives} \\
\hline
\end{tabular}
\caption{Status of four route families. Each obstruction is scoped to the stated
class of arguments.}
\label{tab:routes}
\end{table}

\section{Obstruction from the integrated first dynamic cumulant}
\label{sec:cumulant}

In this section the electronic part of \eqref{eq:H} is the free tight-binding
operator with an arbitrary real static field, and the phonon-mediated interaction is
kept in its exact dynamic form. The route examined here conditions on every static
site field $W$, uses the positive static free-fermion trace as reference, and then
expands the Gaussian-integrated nonzero-Matsubara part of the interaction. Because
the static reference is strictly positive, this route survives all determinant
worries; the question is whether the first coefficient of the expansion is
uniformly small at a linear temperature scale.

\subsection{Setup and the exact first coefficient}

Fix $\tau=\beta^{-1}$ and $\Omega>0$, set $\nu_n=2\pi n/\beta$ and let
$h_W=-tA+\diag(W)$ be the static one-particle matrix, with $A$ the adjacency matrix of
$\Lambda_L$. Two identical, independent spin sectors are present at fixed $W$. For a
mean-zero real Gaussian field $v$ with the thermal covariance of the dynamic part of
the interaction,
\begin{equation}
D_{\mathrm{dyn}}(u)=\frac{U}{\beta}\sum_{n\ne0}r_n e^{i\nu_nu},\qquad
r_n=\frac{\Omega^2}{\nu_n^2+\Omega^2},
\label{eq:cov}
\end{equation}
define the two-spin partition-function ratio obtained by multiplying the
Gaussian-integrated nonzero-mode interaction by $s$,
\begin{equation}
R_W(s)=\frac{\E_v\Tr_{\mathrm{Fock}}\,T
\exp\!\left[-\int_0^\beta\!\left(H_W+\sqrt{s}\sum_xv_x(u)n_x\right)du\right]}
{\Tr_{\mathrm{Fock}}e^{-\beta H_W}} .
\label{eq:RW}
\end{equation}
The denominator is strictly positive and $R_W(0)=1$, so the first coefficient of
$R_W$ equals the first coefficient of $\log R_W$. Both are written
$\sum_xa_x(W)$. Using the force-scaled field $v_x=F_xX_x$, the covariance
\eqref{eq:cov} is exactly the thermal covariance of that field, with
$U_x=F_x^2/K_x$; replacing $n_x$ by $n_x-1$ changes nothing because $D_{\mathrm{dyn}}$
has zero time integral.

\begin{lemma}[Exact first coefficient, no truncation]
\label{lem:first}
For every finite volume, every real static field $W$, every $t,\beta,\Omega>0$ and
every site $x$,
\begin{equation}
a_x(W)=2U\sum_{n\ge1}r_n\sum_{a,b}p_a(x)p_b(x)\,
\frac{\bigl[f(e_a)-f(e_b)\bigr](e_b-e_a)}
{\nu_n^2+(e_b-e_a)^2},
\label{eq:first}
\end{equation}
where $h_W\psi_a=e_a\psi_a$, $p_a(x)=|\psi_a(x)|^2$, the sum over $(a,b)$ is over
ordered eigenstate pairs, and $f(e)=(1+e^{\beta e})^{-1}$. Every summand in
\eqref{eq:first} is nonnegative and diagonal pairs contribute zero.
\end{lemma}

\begin{proof}
Expand the exponential in \eqref{eq:RW} to second order in $\sqrt{s}$ and apply
Gaussian Wick, which is legitimate at finite volume because $\int|v|$ has Gaussian
exponential moments and the electronic traces are finite dimensional. The
second-order term of $R_W$ is
\begin{equation}
\frac12\sum_x\int_0^\beta\!\!\int_0^\beta du\,du'\,
D_{\mathrm{dyn}}(u-u')\bigl\langle T\,n_x(u)n_x(u')\bigr\rangle_W .
\label{eq:second}
\end{equation}
The disconnected part $\langle n_x\rangle\langle n_x\rangle$ is independent of both
times and vanishes because $D_{\mathrm{dyn}}$ has zero integral. The two spin sectors
are independent, so cross-spin connected terms vanish and the two same-spin terms
cancel the prefactor $1/2$. For one spin, the connected density correlation in the
static reference is
\begin{equation}
K_x(u)=\sum_{a,b}p_a(x)p_b(x)f(e_a)\bigl[1-f(e_b)\bigr]e^{u(e_a-e_b)},\qquad
0<u<\beta .
\label{eq:Kx}
\end{equation}
Its spectral representation
$\chi_x(i\nu)=\sum_{a,b}p_ap_b[f(e_a)-f(e_b)](e_b-e_a)/[\nu^2+(e_b-e_a)^2]$ follows
by direct time integration and by pairing $(a,b)$ with $(b,a)$; at nonzero integer
frequencies $\int_0^\beta e^{i\nu(u-u')}du\,du'$ collapses the double time integral
onto $\beta\delta_{\nu_k,\nu_m}$, which cancels the $\beta^{-1}$ in \eqref{eq:cov} and
leaves the factor $U$. Summing over both signs of the bosonic frequency, using
$\chi_x(i\nu)=\chi_x(-i\nu)$, gives \eqref{eq:first}. Positivity is the statement
that $f$ is decreasing: for $e_a<e_b$ the numerator
$[f(e_a)-f(e_b)](e_b-e_a)$ is a positive square divided by a positive difference, and
for equal energies it vanishes.
\end{proof}

\begin{lemma}[Auxiliary bounds]
\label{lem:aux}
Set $q(a)=\sum_{n\ne0}a^2/[(2\pi n)^2+a^2]$ and let $d_x$ be the degree of $x$ in
$\Lambda_L$. Then for every real $W$ and every $\Omega>0$:
\begin{enumerate}[label=(\roman*),leftmargin=2.2em]
\item $q(a)\le a^2/12$, $\chi_x(i\nu)\le\beta/4$, and
$a_x(W)\le U\Omega^2\beta^3/48$;
\item $\sum_{a,b}p_a(x)p_b(x)(e_b-e_a)^2=2t^2d_x$ and
$a_x(W)\le d_xt^2U\Omega^2\beta^5/1440$;
\item for distinct $x,y$ and $\alpha=(1+16\beta t)^{-1}$,
$\bigl|\langle\Delta_x^\dagger\Delta_y\rangle_W\bigr|
\le(25/16)e^{-2\alpha\,d(x,y)}$.
\end{enumerate}
\end{lemma}

\begin{proof}
(i) Since $a^2/[(2\pi n)^2+a^2]\le a^2/(4\pi^2n^2)$ and $\sum_{n\ne0}n^{-2}=\pi^2/3$,
we get $q(a)\le a^2/12$. For the second claim, write $z=e_b-e_a$ and use the mean
value theorem, $(f(e_a)-f(e_b))/(e_b-e_a)=-f'(\xi)\le\beta/4$ together with
$z^2/(\nu^2+z^2)\le1$ and $\sum_{a,b}p_ap_b=1$. Then
$a_x\le(\beta U/4)q(\beta\Omega)\le U\Omega^2\beta^3/48$.

(ii) The spectral variance identity holds because the left side is the variance of
$h_W$ in the site state $|x\rangle$: it equals
$2[(h_W^2)_{xx}-(h_W)_{xx}^2]=2\sum_{y}|(h_W)_{xy}|^2=2t^2d_x$. Since
$z^2/(\nu^2+z^2)\le z^2/\nu^2$, we get $\chi_x(i\nu)\le\beta d_xt^2/(2\nu^2)$, and
$\sum_{n\ne0}\nu_n^{-4}=\beta^4/720$ together with $r_n\le1$ gives the bound.

(iii) Conjugate the hopping by the diagonal multiplication operator
$\exp(\alpha\,d(\,\cdot\,,\text{anchor}))$. The conjugated hopping differs from $h_W$
by a perturbation $P$ with Schur norm at most $4t(e^{\alpha}-1)\le8t\alpha\le
1/(2\beta)$. At each fermionic frequency $\omega_k=(2k+1)\pi/\beta$ the unweighted
inverse of $i\omega_k-h_W$ has norm at most $1/|\omega_k|$ by the spectral theorem,
and the difference of the two inverses is at most $2\|P\|/|\omega_k|^2$, so the
difference series converges absolutely because the identity jumps cancel. With
$\sum_k\omega_k^{-2}=\beta^2/4$ the static slice norm after conjugation is at most
$1+\beta\|P\|/2\le5/4$, and Wick factorisation of the static pair correlator,
$\langle\Delta_x^\dagger\Delta_y\rangle_W=f(h_W)_{xy}^2$, gives
\eqref{eq:clustering} with the stated constant after undoing the conjugation.
\end{proof}

\begin{remark}
The static reference is therefore not where the difficulty lies: item (iii) gives
uniform exponential decay of the static pair correlator at every fixed positive
temperature and for every real field, so any annealed average over static fields
inherits it. Bounds (i) and (ii) also show that the exact first cumulant is always of
order at most $U\Omega^2\beta^3/48$; the next theorem shows that this worst-case
order is attained, at fixed hopping, for a field profile that is not exotic.
\end{remark}

\subsection{The three-site witness}

\begin{theorem}[Fixed-hopping resonance]
\label{thm:witness}
Use the open three-site path $\{$left, middle, right$\}$ with hopping $-t$ and real
fields
\begin{equation}
W_{\mathrm{left}}=W_{\mathrm{right}}=\tfrac{\delta}{2},\qquad
W_{\mathrm{middle}}=-\tfrac{\delta}{2}+\frac{2t^2}{\delta},\qquad \delta>0 .
\label{eq:fields3}
\end{equation}
For every $t,U>0$ and every $0<\Omega\le\tau\le t/4$,
\begin{equation}
a_{\mathrm{left}}(W)\ \ge\ \frac{U\Omega^2}{1200\,\tau^3},
\label{eq:witness}
\end{equation}
and the same holds at the right endpoint.
\end{theorem}

\begin{proof}
The one-particle matrix in the site basis takes the values
$h_3=\tfrac{\delta}{2}$ at the two end sites,
$-\tfrac{\delta}{2}+2t^2/\delta$ at the middle site, and $-t$ on both bonds. Its
antisymmetric vector $(1,0,-1)/\sqrt2$ is an eigenvector with energy $e_+=\delta/2$
and endpoint weight $p_+(\mathrm{left})=1/2$. On
$\mathrm{span}\{(1,0,1)/\sqrt2,(0,1,0)\}$ the remaining block has trace
$2t^2/\delta$ and determinant $-\delta^2/4-t^2$, hence eigenvalues
$e_-=-\delta/2$ and $e_{\mathrm{high}}=\delta/2+2t^2/\delta$; the amplitude ratio of
the middle component to the symmetric endpoint component at energy $e_-$ is
$\delta/(\sqrt2\,t)$, so the endpoint weight is
\begin{equation}
p_-(\mathrm{left})=\frac{A}{2},\qquad
A=\left(1+\frac{\delta^2}{2t^2}\right)^{-1}.
\label{eq:Aweight}
\end{equation}
The eigenvalue separation of the two low states is exactly $\delta$, arbitrarily
small at fixed bare hopping, because increasing the middle barrier $2t^2/\delta$
pays for it.

Keep in \eqref{eq:first} only $n=1$ and the two ordered low-state pairs
$(-,+)$ and $(+,-)$. Their Fermi difference is $\tanh(\beta\delta/4)$ and their
ordered weight products sum to $A/2$, so every retained term is nonnegative and
\begin{equation}
a_{\mathrm{left}}\ge\frac{UA\,\Omega^2\,\delta\tanh(\beta\delta/4)}
{(4\pi^2\tau^2+\Omega^2)\,(4\pi^2\tau^2+\delta^2)} .
\label{eq:witnessgeneral}
\end{equation}
Now set $\delta=4\tau$. For $\Omega\le\tau\le t/4$ we have $\delta\le t$, hence
$A\ge2/3$, and $\tanh(1)>3/4$, and $\pi^2<10$, so the right side of
\eqref{eq:witnessgeneral} is bounded below by
\begin{equation}
\frac{(2/3)\cdot4\cdot(3/4)}{(40+1)(40+16)}\cdot\frac{U\Omega^2}{\tau^3}
=\frac{1}{1148}\cdot\frac{U\Omega^2}{\tau^3}
>\frac{1}{1200}\cdot\frac{U\Omega^2}{\tau^3}.
\label{eq:constant}
\end{equation}
The constants used are elementary and exact: the identity
$\int_0^1x^4(1-x)^4/(1+x^2)dx=22/7-\pi$ with a positive integrand gives $\pi<22/7$,
hence $\pi^2<(22/7)^2<10$, and the Taylor terms of $e^2$ through degree four sum to
$7$ with a strictly positive remainder, giving $\tanh(1)>3/4$. The arithmetic
$(2/3)\cdot4\cdot(3/4)/(41\cdot56)=1/1148>1/1200$ is exact in rational arithmetic.
\end{proof}

\begin{theorem}[Embedding on every square torus]
\label{thm:torus}
Fix $L\ge4$, $t,U>0$, $0<\Omega\le\tau\le t/4$. Let the three collinear torus
sites $(0,0),(0,1),(0,2)$ carry the fields $\delta=4\tau$ of \eqref{eq:fields3} with
$\delta/2=2\tau$ and middle field $t^2/(2\tau)-2\tau$, and let every other site carry
the same real field $B>0$. The bare hopping is unchanged on every torus edge,
including edges into the barrier region, and no bond is deleted. Then there is a
finite $B_0=B_0(L,t,\tau,\Omega)$ such that for every $B\ge B_0$
\begin{equation}
a_{(0,0)}(W_B)\ \ge\ \frac{U\Omega^2}{2400\,\tau^3}.
\label{eq:torus}
\end{equation}
\end{theorem}

\begin{proof}
The three chosen sites induce the open path, including at $L=4$, since the two
endpoints $(0,0)$ and $(0,2)$ are at graph distance $2$. Write the static matrix in
path/complement block form
$h_B=\begin{pmatrix}h_3&V\\ V^*&BI-tA_{\mathrm{c}}\end{pmatrix}$, where
$A_{\mathrm{c}}$ is the adjacency matrix of the complement of the three sites. The
complement spectrum is at least $B-4t$, and $\|V\|\le\sqrt6\,t<4t$. For fixed finite
$L,t,\tau,\Omega$ the low-sector Schur complement
\begin{equation}
S(B,z)=h_3-z-V\bigl(BI-tA_{\mathrm{c}}-z\bigr)^{-1}V^*
\label{eq:schur}
\end{equation}
converges uniformly to $h_3-z$ on any compact set of $z$ away from the three simple
eigenvalues of $h_3$. The inverse of \eqref{eq:schur} is the upper-left $3\times3$
block $[(h_B-z)^{-1}]_{11}$ of the full resolvent, and the remaining blocks
$12$, $21$ and $22$ are $O(B^{-1})$ uniformly on small contours around those
eigenvalues, so the full resolvent converges to
$(h_3-z)^{-1}\oplus0$ there. Contour integration then gives convergence of the
spectral projections, and residues give convergence of the two low eigenvalues and
their local weights. No limit is taken in $L$.

Because every summand of \eqref{eq:first} is nonnegative, retaining only $n=1$ and
the two converging low-state pairs keeps a sum that converges to the right side of
\eqref{eq:witnessgeneral}, which by \eqref{eq:constant} exceeds
$U\Omega^2/(1200\tau^3)>U\Omega^2/(2400\tau^3)$. Hence there is a finite $B_0$ with
\eqref{eq:torus} for all $B\ge B_0$. No explicit value or uniform estimate for $B_0$
is claimed; the prefactor $1/2400$ only absorbs the loss in the embedding.
\end{proof}

\begin{corollary}
\label{cor:scale}
Under the hypotheses of Theorem~\ref{thm:torus}, the exact first local coefficient at
the endpoint satisfies
\begin{equation}
a_{(0,0)}(W_B)\ \ge\ \frac{U}{2400\,C^3\,\Omega}
\qquad\text{whenever}\qquad \tau=C\,\Omega,\quad C\ge1,\quad \Omega\le t/(4C).
\label{eq:linear}
\end{equation}
In particular it diverges as $\Omega\downarrow0$ at every fixed $C$, while
Lemma~\ref{lem:aux}(i) shows it is never larger than $U\Omega^2\beta^3/48$, so the
witness attains the worst-case order to within a factor $50$.
\end{corollary}

\begin{proof}
Substitute $\tau=C\Omega$ in \eqref{eq:torus}. The upper comparison is
Lemma~\ref{lem:aux}(i) evaluated at the same parameters.
\end{proof}

\subsection{Exactly what is obstructed}

\begin{remark}[Obstructed proof class]
\label{rem:class}
The three results above block the following explicitly described class of arguments:
(a) condition on an arbitrary real static field $W$ and use the positive static
free-fermion state of $H_W$ as reference; (b) expand the Gaussian-integrated
nonzero-mode density interaction about that reference, without absorbing or
resumming its connected contributions into a new reference; and (c) require a
uniform absolute local activity, tending to zero at $\tau=C\Omega$, which dominates
the exact local first coefficient \eqref{eq:first} or its per-volume sum. Item (c)
fails, and it fails at fixed bare hopping and at a torus side $L=4$ already, so the
obstruction is not a large-volume or large-$t$ artefact.
\end{remark}

\begin{remark}[What the class argument does not establish]
\label{rem:notclass}
A large first coefficient does not prevent a Taylor or cluster series from
converging, and a large source-free cumulant does not obstruct a source-connected
estimate: resummation can absorb connected contributions into a new reference. The
witness therefore does not prove that every conditional expansion fails. It also does
not obstruct the annealed route, in which the dynamic Gaussian is weighted by the
full Gibbs factor: the field profile of \eqref{eq:fields3} has a middle field
$t^2/(2\tau)-2\tau$ growing like $\tau^{-1}$ at fixed $t$, which is strongly
suppressed under any Gaussian static prior. Finally, no statement is made about
$\Tpair$ itself. The same caveats hold for the comparison with the earlier
two-thirds-power analysis \cite{linear4,linear5}: the dimensionless cubic scale
$U\Omega^2/\tau^3$ was already visible there, and what is new here is the fixed-hopping
witness that attains it after exact Hartree cancellation and full Matsubara
integration, without a determinant zero and without an atomic reference.
\end{remark}

\section{Obstruction from positive Fermion determinants}
\label{sec:det}

A natural way to control the linear scale is to bound, history by history, the
normalised single-particle propagator that multiplies each phonon configuration, and
then to use the positivity of the determinant for the spinful problem to keep the
weights positive. This section shows that positivity and a strict lower bound on the
determinant are together insufficient.

\subsection{The oscillator sum rule, as a classical lemma}

The following standard sum rule is quoted as a lemma because it delimits what a
second-moment argument can do. It is the oscillator $f$-sum rule in imaginary
frequency, and it holds unchanged in the interacting problem because the electronic
operators commute with the oscillator coordinate.

\begin{lemma}[Interacting nonzero-Matsubara sum rule]
\label{lem:fsum}
Let $x$ be any site of \eqref{eq:H} with force-scaled displacement
$\psi_x=F_xX_x$, and let
$C_x(i\nu_k)=\int_0^\beta e^{i\nu_ku}\langle\psi_x(u)\psi_x(0)\rangle du$ be taken in
the exact interacting Gibbs state. Then for every $k\ne0$
\begin{equation}
0\le C_x(i\nu_k)\le\frac{U_x\Omega_x^2}{\nu_k^2}.
\label{eq:fsum}
\end{equation}
In a translation-invariant thermal path representation with normalised Fourier
coefficients $\psi_{x,k}=\beta^{-1}\int_0^\beta\psi_x(u)e^{-i\nu_ku}du$,
\begin{equation}
\E|\psi_{x,k}|^2\le\frac{U_x\Omega_x^2}{4\pi^2k^2\tau},\qquad
\sum_{k\ge1}\E|\psi_{x,k}|^2\le\frac{U_x\Omega_x^2}{24\,\tau},
\label{eq:fsum2}
\end{equation}
and the constant is $1/12$ when both signs of $k$ are summed.
\end{lemma}

\begin{proof}
The oscillator coordinate commutes with the electronic operators and the other
oscillators, so \eqref{eq:dcomm} gives $[H_\Lambda,\psi_x]=-iF_xP_x/M_x$ and
$[\psi_x,[H_\Lambda,\psi_x]]=U_x\Omega_x^2$; the opposite order of the double
commutator has the opposite sign. Let $E_a$ enumerate the exact energies with weights
$p_a=e^{-\beta E_a}/Z$. At nonzero $\nu$ the Lehmann expansion gives
\begin{equation}
C_x(i\nu)=\sum_{a,b}\frac{(p_a-p_b)(E_b-E_a)}{(E_b-E_a)^2+\nu^2}
\bigl|(\psi_x)_{ab}\bigr|^2 .
\label{eq:lehmann}
\end{equation}
Degenerate pairs contribute zero at nonzero $\nu$ and the imaginary parts cancel
between $(a,b)$ and $(b,a)$, so every numerator is nonnegative, and
$\sum_{a,b}(p_a-p_b)(E_b-E_a)|(\psi_x)_{ab}|^2
=\langle[\psi_x,[H_\Lambda,\psi_x]]\rangle=U_x\Omega_x^2$. Replacing each denominator
by $\nu^2$ proves \eqref{eq:fsum}. Then
$\E|\psi_{x,k}|^2=C_x(i\nu_k)/\beta$ in a translation-invariant path measure, and
$\sum_{k\ge1}k^{-2}=\pi^2/6$ gives \eqref{eq:fsum2}. The unbounded oscillator needs no
Fock cutoff. The linear coordinate coupling is relatively form bounded with
arbitrarily small relative bound against the harmonic potential, so
$D(H_\Lambda^{1/2})=D(N_{\mathrm{osc}}^{1/2})$ and $\psi_x|\psi_a\rangle\in
D(H_\Lambda^{1/2})$ for every eigenvector $|\psi_a\rangle$ of $H_\Lambda$. Parseval
then gives
$\langle\psi_a|[\psi_x,[H_\Lambda,\psi_x]]|\psi_a\rangle
=2\sum_b(E_b-E_a)|(\psi_x)_{ab}|^2$ by monotone convergence of the nonnegative
summands, so the double-commutator value $U_x\Omega_x^2$ used above is finite and
exact.
\end{proof}

\begin{remark}
Lemma~\ref{lem:fsum} rules out one specific claim, namely that an electronic
determinant tilt must destroy every nonzero-frequency second-moment bound: the exact
interacting bound has the same high-frequency coefficient as a free oscillator. It
supplies nothing else. It gives no control of the zero mode $k=0$, no Gaussian
concentration, no independence between sites or frequencies, no exponential moment
and no joint rare-event bound. A root-mean-square scale of order
$\sqrt{U_x\Omega_x^2/\tau}$ is not a norm bound on every history, and replacing such a
moment by a uniform bound is exactly the invalid step that
Theorem~\ref{thm:amp} isolates.
\end{remark}

\subsection{Positive determinant, unbounded propagator}

\begin{theorem}[Arbitrary conditional amplification]
\label{thm:amp}
Consider one electronic level coupled to one bath orbital with one-particle matrix
\begin{equation}
h(u)=\begin{pmatrix}\varepsilon(u)&-v\\ -v&0\end{pmatrix},\qquad v>0,\qquad
\varepsilon(u)=\begin{cases}-B,&0<u<1,\\ +B,&1<u<2,\end{cases}
\label{eq:impurity}
\end{equation}
at $\beta=2$ and $B>0$. Let $U_1=e^{-h_1}$ and $U_2=e^{-h_2}$ be the propagators of
the two slices, $U=U_2U_1$ the monodromy, and
\begin{equation}
\Ggreen(1,0)=-U_1\left(1+U_2U_1\right)^{-1}
=-\left(U_1^{-1}+U_2\right)^{-1}
\label{eq:Gdef}
\end{equation}
the antiperiodic two-point Green matrix at the slice boundary. Then
$\det(1+U)>1$, every eigenvalue of $U$ is positive, and
\begin{equation}
\bigl|\Ggreen(1,0)_{00}\bigr|\longrightarrow\infty\quad(B\to\infty),\qquad
\bigl|\Ggreen(1,0)_{00}\bigr|\sim\frac{B^2}{2v^2}\quad(B\to\infty).
\label{eq:amp}
\end{equation}
\end{theorem}

\begin{proof}
Both $U_1$ and $U_2$ are symmetric and positive definite. Their product is similar to
the symmetric positive-definite matrix $U_1^{1/2}U_2U_1^{1/2}$, so all eigenvalues
$\lambda_j$ of $U$ are positive, and
$\det U=e^{-\Tr h_1-\Tr h_2}=1$ because the two slices carry opposite level shifts
and the bath level is zero. Hence $\det(1+U)=\prod_j(1+\lambda_j)>1$.

Set $d=\sqrt{B^2/4+v^2}$, $a=B/(2d)$, $c=\cosh d$, $s=\sinh d$. The two off-diagonal
entries cancel in $U_1^{-1}+U_2$, giving
$U_1^{-1}+U_2=2e^{-B/2}\diag(c-as,\,c+as)$ with both diagonal entries positive
because $0<a<1$. Therefore
\begin{equation}
\bigl|\Ggreen(1,0)_{00}\bigr|
=\frac{e^{B/2}}{(1-a)e^{d}+(1+a)e^{-d}}
=\frac{1}{(1-a)e^{d-B/2}+(1+a)e^{-d-B/2}} .
\label{eq:Gclosed}
\end{equation}
Since $d\ge B/2$ and $d+B/2\ge B$ we have $d(d+B/2)\ge B^2/2$, hence
$1-a=v^2/[d(d+B/2)]\le2v^2/B^2$, and since
$d-B/2=v^2/(d+B/2)\le v^2/B$ and $1+a\le2$,
\begin{equation}
\bigl|\Ggreen(1,0)_{00}\bigr|\ \ge\
\frac{1}{(2v^2/B^2)\,e^{v^2/B}+2e^{-B}},
\label{eq:Glower}
\end{equation}
whose denominator tends to zero as $B\to\infty$. This proves divergence without any
floating-point asymptotic fit. For the rate, expand
$d-B/2=v^2/B+O(B^{-3})$ and $1-a=2v^2/B^2+O(B^{-4})$ in \eqref{eq:Gclosed}.
\end{proof}

\begin{remark}[Smooth and finite-action histories]
\label{rem:smooth}
The step profile of \eqref{eq:impurity} is not a finite-action harmonic-oscillator
path. This does not weaken the obstruction. Approximate the step in $L^1$ by smooth
periodic histories, closing the transitions near $u=0,1,2$; the finite-dimensional
evolution depends continuously on the history through its integral equation and
Gronwall's inequality, and the determinant and the Green matrix are continuous
wherever $\det(1+U)\ne0$. The strict determinant and amplification inequalities
therefore persist for all sufficiently close smooth histories, and for each of them
the oscillator action $\int[(\varepsilon')^2/\Omega^2+\varepsilon^2]/(2U)$ is finite for every
$U,\Omega>0$, though possibly large. No lower bound on the probability of those
histories and no divergence of an averaged propagator is asserted: a
determinant-weighted integral may suppress them. What is obstructed is a uniform
pathwise bound obtained from determinant positivity alone, not an annealed estimate
that retains the weights.
\end{remark}

\begin{corollary}[Interval-certified two-slice fixture]
\label{cor:fixture}
At $B=6$ and $v=1/2$ in \eqref{eq:impurity},
\begin{equation}
\det(1+U)>15,\qquad \bigl|\Ggreen(1,0)_{00}\bigr|>52.89,
\qquad \bigl|\Ggreen(1,0)_{00}\bigr|=52.8957540935\ldots
\label{eq:fixture}
\end{equation}
The numbers were obtained from the exact closed form \eqref{eq:Gclosed} in
outward-rounded interval arithmetic, evaluated with two independently written
implementations and unlike precisions, so each enclosure provably contains the true
value. They are calculated interval fixtures, not measured data. The fixture alone
rejects the proposed uniform bound $|\Ggreen(1,0)_{00}|\le1$; the asymptotic theorem
is analytical and is not established by the fixture.
\end{corollary}

\subsection{Supporting algebra: Ward identities, and two dimer facts}

\begin{remark}[Finite-matrix Ward and Schur algebra]
\label{rem:ward}
Split a finite Matsubara matrix into positive and negative frequency blocks as
$A=\begin{pmatrix}R_++iD_+&-V\\ -V^*&R_- -iD_-\end{pmatrix}$ with
$R_\pm=R_\pm^*$, $D_\pm>0$ and $G=A^{-1}$. Each diagonal block is invertible and
\begin{equation}
\bigl\|D_\pm^{1/2}A_\pm^{-1}D_\pm^{1/2}\bigr\|\le1,
\label{eq:ward1}
\end{equation}
because the normalised block is $H\pm i$ with $H$ Hermitian and its inverse singular
values are $1/\sqrt{1+\lambda(H)^2}$. With
$Q=VA_-^{-1}D_-(A_-^{-1})^*V^*$ one has $0\le Q\le VD_-^{-1}V^*=:\tilde Q$, and Schur
complementation together with $\operatorname{Im}(S^{-1})=-S^{-1}(\operatorname{Im}S)
(S^{-1})^*$ gives
\begin{equation}
G_{+-}D_-G_{+-}^*=G_{++}QG_{++}^*,\qquad
-\operatorname{Im}G_{++}=G_{++}(D_+-Q)G_{++}^* .
\label{eq:ward2}
\end{equation}
If $\alpha=\|D_+^{-1/2}VD_-^{-1/2}\|^2<1$ then $Q\le\alpha D_+$ and
\begin{equation}
\left(G_{+-}D_-G_{+-}^*\right)_{nn}\le\frac{\alpha}{1-\alpha}
\left(-\operatorname{Im}G_{++}\right)_{nn},
\label{eq:ward3}
\end{equation}
and the threshold $\alpha=1$ is sharp for positive dissipation in this abstract
algebra, since scalar blocks $R_\pm=0$, $|V|^2=\alpha D_+D_-$ give $Q=\alpha D_+$.
For diagonal $D_\pm=\diag(D_{\pm,m})$ the row sums satisfy exactly
\begin{equation}
\sum_m|(G_{++})_{nm}|^2D_{+,m}+\sum_m|(G_{+-})_{nm}|^2D_{-,m}
=-\operatorname{Im}G_{nn}+2\sum_m|(G_{+-})_{nm}|^2D_{-,m},
\label{eq:ward4}
\end{equation}
where the factor two is not optional. These are finite-matrix identities under the
stated hypotheses; passing to an infinite Matsubara operator requires convergence and
domain control that the identities do not supply, and \eqref{eq:ward3} cannot be
averaged by replacing $\alpha$ with $\E\alpha$, since histories with $\alpha\ge1$ are
not excluded by a second moment and $\alpha$, $G$ and the determinant weight are
correlated. The unconditional sum rule of Lemma~\ref{lem:fsum} does give the
finite-grid moment estimate
$\E\alpha\le U\Omega^2(4\pi^4\tau^3)^{-1}S$ with
$S=\sum_{k\ge1}H_{\mathrm{odd}}(k)/k^3$ and
$H_{\mathrm{odd}}(k)=\sum_{j=0}^{k-1}(2j+1)^{-1}\le1+\tfrac12\log(2k-1)$, hence
$S\le C_{\mathrm{aux}}$ with $67/50<C_{\mathrm{aux}}<27/20$, an enclosure that is
certified in outward-rounded rational arithmetic; the $\tau^{-3}$ scale
reproduces a two-thirds-power small-moment scale and is not a linear temperature law.
\end{remark}

\begin{remark}[Two static dimer anchors]
\label{rem:dimer}
In the two-site spinful Holstein model at fixed total electron number $N$, with
$a_\pm=(a_1\pm a_2)/\sqrt2$ and $G^2=U\Omega/2$ in the notation of \eqref{eq:U}, the
centre-of-mass oscillator decouples exactly after a shift,
\begin{equation}
\Omega a_+^*a_+-\frac{GN}{\sqrt2}(a_++a_+^*)=\Omega c_+^*c_+-\frac{G^2N^2}{2\Omega},
\qquad c_+=a_+-\frac{GN}{\sqrt2\,\Omega},
\label{eq:com}
\end{equation}
while the unshifted centre-of-mass number does not generally commute with $H$. In the
two-electron singlet sector at a frozen relative field $q$ the electronic matrix
$h(q)$ has entries $-q$, $q$, $0$ on the diagonal and $-\sqrt2\,t$ in the off-diagonal
pair, with elastic energy $q^2/(4U)$. The lower Born-Oppenheimer surface
$q^2/(4U)-\sqrt{q^2+4t^2}$ has minima at $q=\pm2\sqrt{U^2-t^2}$ for $U>t>0$, and at
either minimum the bare static pair expectation is
\begin{equation}
\langle\Delta_1^\dagger\Delta_2\rangle_{\mathrm{el}}=\frac{t^2}{4U^2}
\qquad\left(\frac{1}{64}\ \text{at}\ U=4,\ t=1,\ \text{against}\ \frac{1}{68}
\ \text{at the unrelaxed field}\ q=2U\right).
\label{eq:dimermin}
\end{equation}
Virtual electronic admixture within one static well therefore already contributes a
nonzero short-range bare pair correlator, which obstructs a Peierls argument that
identifies every short-range bare pair correlator with a phonon-assisted inter-well
tunnelling amplitude. It says nothing about long-distance order, and the static
energy landscape is not a many-site dressed-coherence bound.
\end{remark}

\section{What survives: the case of vanishing effective attraction}
\label{sec:survives}

The obstructions above concern the attractive Holstein problem at fixed $U>0$ and the
propagator machinery used to control it. There is one case in which a
frequency-uniform ceiling is available: the pure bond-SSH model \eqref{eq:Hssh}, where
the local attraction vanishes and the only coupling is the Peierls one. This closes
the case $\Ueff=0$ that \cite{sshbond3} left open. All constants below are exact
rationals certified by exact-rational and outward-rounded interval replay.

\subsection{High-temperature clustering}

Fix $L\ge4$, the bond relaxation energy $\Gb=g^2/K$, and
\begin{equation}
c_0=2^{-10},\qquad \xi_0=2^{-9}\beta^{-1},\qquad
\varepsilon=\frac{2c_0}{1-c_0}=\frac{2}{1023},\qquad
q_{\mathrm{GC}}=48\varepsilon=\frac{32}{341},
\label{eq:constants}
\end{equation}
so that $q_{\mathrm{GC}}<1/10$, together with
\begin{equation}
A_{\mathrm{GC}}=\frac{3}{1-q_{\mathrm{GC}}}=\frac{341}{103},\qquad
b_{\mathrm{GC}}=\frac{3q_{\mathrm{GC}}}{1-q_{\mathrm{GC}}}=\frac{32}{103},\qquad
R_{\max}=1+\frac{4A_{\mathrm{GC}}\,q_{\mathrm{GC}}}{(1-q_{\mathrm{GC}})^2}
=\frac{24718511}{9834543}<3 .
\label{eq:constants2}
\end{equation}
The source-augmented graph of \eqref{eq:Hssh} adds $|\Lambda_L|$ source edges
$e'_x=(x,x+r)$ to the physical bonds, raising the maximum vertex degree from $4$ to
$6$, with source vertex operator
$Y_{e'}(\theta)=e^{i\theta}\Delta_x^\dagger\Delta_{x+r}+e^{-i\theta}
\Delta_{x+r}^\dagger\Delta_x$ of operator norm $1$ and spectrum $\{-1,0,1\}$.

\begin{theorem}[Frequency- and $\mu$-uniform pair clustering]
\label{thm:gc}
For every $L\ge4$, every real $t$ and $g$, every $K>0$, every $M>0$, every
$\mu\in\R$, and every $\beta>0$ with
$\beta|t|+\beta\Gb\le2^{-10}$, the equal-time on-site pair correlator satisfies
\begin{equation}
\left|\langle\Delta_x^\dagger\Delta_y\rangle_{\Lambda_L}\right|
\le A_{\mathrm{GC}}\,q_{\mathrm{GC}}^{d(x,y)}
=\frac{341}{103}\left(\frac{32}{341}\right)^{d(x,y)}
\label{eq:gcb}
\end{equation}
for all distinct $x,y\in\Lambda_L$, and
\begin{equation}
\sup_{x\in\Lambda_L}\sum_{y\in\Lambda_L}
\left|\langle\Delta_x^\dagger\Delta_y\rangle_{\Lambda_L}\right|
\le R_{\max}=\frac{24718511}{9834543}<3 .
\label{eq:rmax}
\end{equation}
Moreover, for every fixed nonzero $r$, the density
$p_{\Lambda_L}(\xi)=|\Lambda_L|^{-1}\log Z_{\Lambda_L}(\xi)$ of the
source-augmented model extends to a zero-free analytic function of $\xi$ on
$|\xi|<\xi_0$, and for real $\xi\in[0,\xi_0/2]$ obeys
\begin{equation}
p_{\Lambda_L}(\xi)-p_{\Lambda_L}(0)\le
2\beta A_{\mathrm{GC}}\,q_{\mathrm{GC}}^{d(0,r)}\xi+\frac{4\xi^2}{\xi_0^2}.
\label{eq:source}
\end{equation}
\end{theorem}

\begin{proof}[Proof summary]
Five exact steps, each with a certified constant.
\emph{(1) Analyticity and polymer factorisation.} For complex $\xi$ the semigroup
$-H_\Lambda(\xi)$ is norm analytic and the Dyson series is trace-norm summable, since
the phonon trace factors supply a factor
$(2\sinh(\beta\Omega/4(N+1)))^{-|\mathcal{E}_\Lambda|}$ that grows only polynomially
in the order $N$. Each insertion operator, $T_b$, $Y_{e'}$ and
$\Delta_x^\dagger\Delta_y$, is even in fermions, so operators on vertex-disjoint
supports commute identically and the normalised partition function factorises into a
sum over vertex-disjoint connected polymers $E$.
\emph{(2) Activity bound without an atomic gap.} Summing the multiplicities on a
connected polymer, bounding the electronic trace by the multilinear Schatten-H\"older
inequality, which cancels the single-site weight $z_0^{|S_E|}$ exactly, and evaluating
the exact positive Gaussian Wick contractions of the bond oscillators, each pair of
which contributes the frequency-uniform integral
$\int_0^\beta\int_0^\beta C_\beta(|u-v|)\,du\,dv=\beta/K$ with $C_\beta$ the thermal
position covariance of one oscillator, gives
\begin{equation}
|w(E)|\le\left(e^{2\beta|t|+2\beta\Gb}-1\right)^{|E_{\mathrm{phys}}|}
\left(e^{\beta|\xi|}-1\right)^{|E_{\mathrm{src}}|},\qquad
|w_{xy}(E)|\le\left(e^{2\beta|t|+2\beta\Gb}-1\right)^{|E|},
\label{eq:activity}
\end{equation}
and $e^x-1\le x/(1-x/2)$ for $x\in[0,2)$ turns the two exponents into the common
bound $\varepsilon=2/1023$, using $2\beta|t|+2\beta\Gb\le2^{-9}$ and
$\beta|\xi|\le2^{-9}$.
\emph{(3) Graph counting.} A depth-first traversal rooted at any vertex shows that a
degree-$D$ graph has at most $D^{2k}$ connected edge sets of size $k$ containing a
given vertex, and connectivity gives $|S_E|\le k+1$. With $D=4$ the Kotecký-Preiss
sum is then at most $3\sum_{k\ge1}(48\varepsilon)^k=3q_{\mathrm{GC}}/(1-q_{\mathrm{GC}})
=b_{\mathrm{GC}}<1$, which is the origin of the factor $48$.
\emph{(4) Support deletion and the source derivative.} The Kotecký-Preiss criterion
gives $|\widetilde Z_{\Lambda_L\setminus S}/\widetilde Z_{\Lambda_L}|\le e^{|S|}$,
so the marked polymer sum is bounded by $A_{\mathrm{GC}}q_{\mathrm{GC}}^d$ because any
connected polymer connecting $x$ and $y$ has at least $d(x,y)$ edges. The distance
shell of a torus holds at most $4d$ sites, which gives \eqref{eq:rmax}. For the
source, the degree rises to $6$ and the criterion holds with
$q_{\mathrm{src}}=108\varepsilon=72/341$ and $b_{\mathrm{src}}=216/269<1$; since
translation invariance gives
$p'(0)=2\beta\operatorname{Re}(e^{i\theta}\langle\Delta_0^\dagger\Delta_r\rangle)$, and Cauchy's
formula on $|\xi|=\xi_0$ bounds the remainder of order at least two by
$4\xi^2/\xi_0^2$, \eqref{eq:source} follows.
\emph{(5) Certification.} Every rational constant in \eqref{eq:constants} and
\eqref{eq:constants2} and the integer polynomial identities of
Theorem~\ref{thm:routes}(i) were replayed in exact rational arithmetic and, where an enclosure is
needed, in outward-rounded interval arithmetic.
\end{proof}

\begin{lemma}[Boundary spring-constant bracket]
\label{lem:canon}
Let $b\ge2$, cut $\Lambda_L$ into rectangular blocks $B_\alpha$ of side lengths in
$[b,2b-1]$, let $\partial$ be the set of physical bonds with endpoints in different
blocks and $\partial_{\mathrm{src}}$ the set of source edges with endpoints in
different blocks, and for $K'>0$ define the block-decoupled operator
\[
H_{\mathrm{blk}}(\xi;K')=\sum_\alpha H_{B_\alpha}(\xi)
+\sum_{b\in\partial}\Bigl(\frac{P_b^2}{2M}+\frac{K'}{2}X_b^2\Bigr),
\]
where $H_{B_\alpha}(\xi)$ keeps only the vertices, physical bonds and source edges
internal to $B_\alpha$. Then for every real $\xi$, in the full trace and in every
fixed total-number and fixed spin-number sector,
\begin{equation}
\bigl|\log Z_{\Lambda_L}(\xi)-\log Z_{\mathrm{blk}}(\xi;K)\bigr|
\le C_{\mathrm{bdry}}|\partial|+\beta|\xi|\,|\partial_{\mathrm{src}}|,
\label{eq:bdry}
\end{equation}
with
$C_{\mathrm{bdry}}=\beta\bigl(2|t|+4\Gb\bigr)
+\max\bigl\{\log\frac{Z_{\mathrm{ph}}(K/2)}{Z_{\mathrm{ph}}(K)},
\log\frac{Z_{\mathrm{ph}}(K)}{Z_{\mathrm{ph}}(2K)}\bigr\}$ and
$Z_{\mathrm{ph}}(K')=(2\sinh(\beta\sqrt{K'/M}/2))^{-1}$.
\end{lemma}
\begin{proof}
On a bond $b=\{x,y\}$ the operator $T_b$ is the sum of two spin terms, each of
operator norm $1$ acting on orthogonal spin sectors, and $X_b$ is a phonon operator on
a separate tensor factor, so $\|T_b\|_{\mathrm{op}}=2$ and $[X_b,T_b]=0$. Young's
inequality $2|g||X_b|\le KX_b^2/4+4g^2/K$ then gives, as quadratic forms,
\[
-\Bigl(2|t|+\frac{4g^2}{K}\Bigr)I-\frac{K}{4}X_b^2
\le(-t+gX_b)T_b\le
\Bigl(2|t|+\frac{4g^2}{K}\Bigr)I+\frac{K}{4}X_b^2 .
\]
Adding the boundary oscillator potential $\frac{K}{2}X_b^2$ to both sides gives
$\frac{K}{2}X_b^2-\frac{K}{4}X_b^2=\frac{K/2}{2}X_b^2$ from below and
$\frac{K}{2}X_b^2+\frac{K}{4}X_b^2\le\frac{2K}{2}X_b^2$ from above, while a crossing
source edge contributes $\pm\xi Y_{e'}(\theta)$ with $\|Y_{e'}(\theta)\|_{\mathrm{op}}=1$.
Summing over the crossing bonds yields the two-sided operator inequality
\[
H_{\mathrm{blk}}(\xi;K/2)-|\partial|\Bigl(2|t|+\frac{4g^2}{K}\Bigr)-|\xi||\partial_{\mathrm{src}}|
\le H_{\Lambda_L}(\xi)\le
H_{\mathrm{blk}}(\xi;2K)+|\partial|\Bigl(2|t|+\frac{4g^2}{K}\Bigr)+|\xi||\partial_{\mathrm{src}}|
\]
on the form domain. Every operator here is self-adjoint with compact resolvent in each
sector, so min-max monotonicity of eigenvalues bounds the two traces by
$E=\beta|\partial|(2|t|+4\Gb)+\beta|\xi||\partial_{\mathrm{src}}|$. The boundary
oscillators act on tensor factors decoupled from all blocks and commute with the
number operators, so their partition functions factor out identically in every sector,
$\log Z_{\mathrm{blk}}(\xi;K/2)-\log Z_{\mathrm{blk}}(\xi;K)
=|\partial|\log[Z_{\mathrm{ph}}(K/2)/Z_{\mathrm{ph}}(K)]$ and likewise for $K$ to $2K$;
inserting $\log Z_{\mathrm{blk}}(\xi;K)$ between the two bounds gives \eqref{eq:bdry}.
\end{proof}

\begin{theorem}[Thermodynamic canonical decay]
\label{thm:canon}
Under the same condition $\beta|t|+\beta\Gb\le2^{-10}$, let $L_j\to\infty$ and let
$N_j/L_j^2\to n\in(0,2)$ (or balanced spin sectors with $N_{j\sigma}/L_j^2\to a\in(0,1)$).
Then for every fixed $r\ne0$
\begin{equation}
\limsup_{j\to\infty}\left|\langle\Delta_0^\dagger\Delta_r\rangle_{L_j,\mathrm{can}}
\right|\le A_{\mathrm{GC}}\,q_{\mathrm{GC}}^{\|r\|_1}.
\label{eq:canon}
\end{equation}
Consequently the pair-decay temperature of the pure bond-SSH model obeys
\begin{equation}
k_B\Tpair\le1024\left(|t|+\Gb\right)
\label{eq:tpair}
\end{equation}
uniformly in $\Omega>0$, in $\mu\in\R$, in the interior density $n\in(0,2)$, and in
the torus side.
\end{theorem}

\begin{proof}[Proof summary]
Write $\widehat Z$ for the partition function divided by the free-phonon factor. Cut
$\Lambda_L$ into rectangular blocks of side $[b,2b-1]$ ($b\ge2$) and decouple the
blocks by replacing each boundary spring constant $K$ of a physical bond crossing the
blocks by $K/2$ from below and by $2K$ from above; Young's inequality
$2|g||X_b|\le KX_b^2/4+4g^2/K$ and min-max eigenvalue monotonicity then bracket
$\log Z_{\Lambda_L}(\xi)$ by the two decoupled values with boundary error at most
$C_{\mathrm{bdry}}|\partial|+\beta|\xi||\partial_{\mathrm{src}}|$, where
$C_{\mathrm{bdry}}=\beta(2|t|+4g^2/K)$ plus the logarithm of the ratio of single
oscillator partition functions at the two spring constants, and
$|\partial|\le4|\Lambda_L|/b$ (Lemma~\ref{lem:canon}). Distributing the $N$ particles
among the blocks adds at most $O(b^{-2}\log b)|\Lambda_L|$, so
$\log\widehat Z_{\Lambda_L,N}(0)$ is superadditive up to $O(b^{-1})|\Lambda_L|$ and the
canonical free-energy density $f_{\mathrm{can}}$ is concave in $n$. Its supporting
slope at an interior $n$ is $-\beta\mu$ for an admitted $\mu$, and the finite-volume
Legendre discrepancy
$\delta_L=L^{-2}(\log\widehat Z_{\Lambda_L}(0;\mu)-\log\widehat Z_{\Lambda_L,N}(0)
-\beta\mu N)$ is nonnegative and tends to $0$. Since the Hamiltonian is linear in the
real source $\xi$ and self-adjoint, $\xi\mapsto L^{-2}\log\widehat Z_{\Lambda_L,N}(\xi)$
is convex, so its derivative at $\xi=0$ is bounded by the difference quotient, i.e.\ by
$(p_{\Lambda_L}(\xi)-p_{\Lambda_L}(0))/\xi+\delta_L/\xi$ there; Theorem~\ref{thm:gc}
holds uniformly in $\mu$, so taking $\limsup_{j\to\infty}$ at fixed $\xi$, then
$\xi\downarrow0$, then optimising over the phase $\theta$ gives \eqref{eq:canon}.
Inverting $\beta|t|+\beta\Gb\le2^{-10}$ gives \eqref{eq:tpair}.
\end{proof}

\begin{remark}
Theorem~\ref{thm:canon} is uniform in the frequency, but it is not a linear ceiling:
the constant $1024(|t|+\Gb)$ does not vanish as $\Omega\downarrow0$, so it says
nothing at the scale $\Tlin$. It closes the high-temperature regime of the
$\Ueff=0$ case that the strong-coupling analyses leave open. The Holstein ceiling of
\cite{stage6} and the attractive-Hubbard ceiling of \cite{sshbond3} both carry the rate
$t^2/\Ueff$, which is infinite at $\Ueff=0$, whereas \eqref{eq:tpair} is finite there;
at $t=0$ the threshold $2^{13}\Gb$ of \cite{sshbond3} becomes $2^{10}\Gb$ through
Theorem~\ref{thm:gc}. Neither paper treats the pure Peierls case, so this is a new
regime and not a sharper constant in theirs.
\end{remark}

\subsection{Three further routes obstructed at \texorpdfstring{$\Ueff=0$}{U\_eff = 0}}

\begin{theorem}[Exact route obstructions]
\label{thm:routes}
The following three statements hold for the pure bond-SSH model.
\begin{enumerate}[label=(\roman*),leftmargin=2.2em]
\item \emph{No phonon-induced instantaneous atomic gap.} Complete the square in each
bond oscillator at $t=0$, $\mu=0$:
$H_{\mathrm{ph},b}+gX_bT_b=\frac{P_b^2}{2M}+\frac{K}{2}(X_b+\frac{g}{K}T_b)^2
-\frac{g^2}{2K}T_b^2$, so the antiadiabatic limit is
$H^{\mathrm{inst}}_{\mathrm{eff}}=-\frac{g^2}{2K}\sum_bT_b^2$ exactly. On the
four-site ring $C_4$ at half filling $N=4$ this operator conserves
$(N_\uparrow,N_\downarrow)$ and every local parity $(-1)^{n_x}$, hence $N_{\mathrm{odd}}$.
With $\Sigma_4=\sum_{b\in\mathcal{E}_{C_4}}T_b^2$, one has
$\spec(\Sigma_4|_{N_{\mathrm{odd}}=0})=\{0,4,4,4,8,12\}$ and
$\max\spec(\Sigma_4|_{N_{\mathrm{odd}}>0})=12$ as well, so
$\min\spec(H^{\mathrm{inst}}_{\mathrm{eff}})_{N_{\mathrm{odd}}=0}=
\min\spec(H^{\mathrm{inst}}_{\mathrm{eff}})_{N_{\mathrm{odd}}>0}=-6g^2/K$ and the
phonon-induced gap
$\Delta_{\mathrm{eff}}=0$.
\item \emph{Non-terminating bond Lang-Firsov expansion.} For bonds
$b=\{x,y\}$, $b'=\{y,z\}$ sharing the vertex $y$,
\begin{equation}
[T_b,T_{b'}]=\sum_{\sigma}\left(c_{x\sigma}^\dagger c_{z\sigma}
-c_{z\sigma}^\dagger c_{x\sigma}\right)\ne0,\qquad
\bigl\|[T_b,T_{b'}]\bigr\|_{\mathrm{op}}=2 .
\label{eq:bondcomm}
\end{equation}
Under the bond generator $S=-i\frac{g}{K}\sum_bP_bT_b$,
$\operatorname{ad}_S^k(H_\Lambda)$ contains degree-$k$ momentum products
$(g/K)^kP_{b_1}\cdots P_{b_k}\operatorname{ad}_{T_{b_1}}\cdots
\operatorname{ad}_{T_{b_{k-1}}}(T_{b_k})$ along contiguous bond walks for every
$k\ge2$, whose thermal root-mean-square scale per factor is
$|g|K^{-1}\langle P_b^2\rangle_{0,\beta}^{1/2}
=\sqrt{g^2(2K\Omega)^{-1}\coth(\beta\Omega/2)}$, which diverges as $\Omega\downarrow0$
at fixed $\beta,t,g,K$. At first order,
$[S,V_t]=i(gt/K)\sum_{b,b'}P_b[T_b,T_{b'}]$ and
$[S,V_g]=-\frac{g^2}{K}\sum_bT_b^2-i\frac{g^2}{K}\sum_{b\ne b'}P_bX_{b'}[T_b,T_{b'}]$.
\item \emph{Hopping-uniform exponential clustering is impossible.} At $g=0$,
$\mu=0$, $t>0$ and half filling, $\langle\Delta_0^\dagger\Delta_r\rangle=G(r)^2$ with
$G$ the free-fermion correlator. Along the staircase ray $r_m=(m,m-1)$ of graph
distance $d_m=2m-1$,
\begin{equation}
\lim_{\beta t\to\infty}\lim_{L\to\infty}G(r_m)
=\frac{(-1)^{m-1}2}{\pi^2(2m-1)},\qquad
\lim_{\beta t\to\infty}\lim_{L\to\infty}
\langle\Delta_0^\dagger\Delta_{r_m}\rangle=\frac{4}{\pi^4d_m^2}.
\label{eq:staircase}
\end{equation}
Hence for every $A>0$ and $q\in(0,1)$ there are $m$, $\beta t$ and $L\ge4$ with
$|\langle\Delta_0^\dagger\Delta_{r_m}\rangle|>Aq^{d_m}$: no exponential clustering
bound with hopping-independent constants can follow from a condition on $\beta\Gb$
alone.
\end{enumerate}
\end{theorem}

\begin{proof}
(i) Completing the square is exact, and $T_b^2$ expands as
$h_{b,\uparrow}^2+h_{b,\downarrow}^2+2h_{b,\uparrow}h_{b,\downarrow}$ with
$h_{b,\sigma}^2=(n_{x\sigma}-n_{y\sigma})^2$ diagonal and the cross term transferring
both spins at once, so $(-1)^{n_z}$ and hence $N_{\mathrm{odd}}=\sum_z(n_{z\uparrow}
-n_{z\downarrow})^2$ are conserved. On the paired sector the six hard-core pair states
give the integer matrix whose exact spectrum is $\{0,4,4,4,8,12\}$ with eigenvector
$(1,2,1,1,2,1)^{\mathrm T}$ for the value $12$, certified by the integer polynomial
identity $M(M-4)(M-8)(M-12)=0$. On the fully singly occupied sector
$\frac14T_b^2=\frac12(1-P_{xy})=\frac14-2\mathbf S_x\cdot\mathbf S_y$, twice a singlet
projector, and the four-site Heisenberg ring gives the same spectrum on the
spin-balanced subspace, which is where the value $12$ is attained a second time. On
the remaining invariant sectors the spectrum of $\Sigma_4$ is computed in exact
integer arithmetic as well, and its maximum is strictly smaller: $10$ on the
$24$-dimensional spin-balanced $\{N_{\mathrm{odd}}=2\}$ sector, $8$ on the two singly
occupied $S_z=\pm1$ sectors of dimension $4$, $6$ on the two polarised
$\{N_{\mathrm{odd}}=2\}$ sectors of dimension $12$, and $0$ on the two fully polarised
sectors. The half-filled space has $\binom84=70$ states, and the verification covers
all $70$: the spectrum is computed on nine individually invariant sectors, which pair
into six classes under the global spin flip. Hence the maximum of $\Sigma_4$ equals
$12$ on the unpaired sectors as well as on the paired sector.
(ii) For each spin species, canonical anticommutation gives
$[h_{b,\sigma},h_{b',\sigma}]=c_{x\sigma}^\dagger c_{z\sigma}-c_{z\sigma}^\dagger
c_{x\sigma}$, and the two spin species commute, so the square of $i$ times the
commutator is $(n_{x\sigma}-n_{z\sigma})^2$ of norm one per spin and the operator norm
is $2$. Because $[P_b,P_{b'}]$ vanishes and $[P_b,X_{b'}]=-i\delta_{bb'}$, each
application of $\operatorname{ad}_S$ appends one momentum factor and one nested commutator
of hopping operators, which is a nonvanishing end-to-end current of norm $2$ along any
path of distinct bonds; no finite order terminates the series. The single-oscillator
moment is $\langle P_b^2\rangle_{0,\beta}=\frac{K}{2\Omega}\coth(\beta\Omega/2)$.
(iii) At $g=0$ the two spin species are independent, so the singlet-pair correlator
factors into a square of the single-particle correlator, and Fourier inversion gives
$G_{\Lambda_L,\beta}(r)$ as the stated lattice sum. Taking $L\to\infty$ turns the sum
into a Riemann integral, and $\beta t\to\infty$ restricts it to the Fermi sea
$\{\cos k_x+\cos k_y>0\}$ by dominated convergence because the Fermi surface has
measure zero. In rotated coordinates $u=(k_x+k_y)/2$, $v=(k_x-k_y)/2$ the sea is the
diamond $|u|<\pi/2$, $|v|<\pi/2$ up to a Jacobian $2$, and the two separated integrals
give \eqref{eq:staircase}. Since $\pi^{-4}d^{-2}$ dominates $Aq^{d}$ for any $q<1$ at
large $d$, the negative statement follows. As an outward-rounded rational check, at
$m=2$ one has $4/(9\pi^4)\approx4.5627\times10^{-3}>A_{\mathrm{GC}}q_{\mathrm{GC}}^3
\approx2.7359\times10^{-3}$.
\end{proof}

\begin{remark}
Item (i) says that the antiadiabatic bond model has no instantaneous two-body operator
separating paired from unpaired sectors on the smallest allowed torus, so an argument
that first takes $\Omega\to\infty$ and then uses an atomic gap cannot be correct.
Item (ii) says that a bond Lang-Firsov transformation is not a finite-order device for
Peierls coupling, in contrast with the site coupling of \eqref{eq:H} where the
transformation acts diagonally. Item (iii) is a statement about what a condition on
$\beta\Gb$ alone can imply: the algebraic free-fermion tail is invisible to the
one-scale polymer estimate when the hopping is large, which is why
Theorem~\ref{thm:gc} requires the combined condition
$\beta|t|+\beta\Gb\le2^{-10}$ rather than a condition on $\beta\Gb$ alone.
\end{remark}

\section{What is not claimed}
\label{sec:notclaimed}

The results of this note are model statements at the level of partition functions,
correlators and Hamiltonians. They are not measurements, not predictions for a
material, and not claims about a measured critical or ordering temperature. In
particular:

\begin{enumerate}[label=(\alph*),leftmargin=2.2em]
\item No upper bound on $\Tpair$ for the attractive Holstein model, or for any model
with $\Ueff>0$, is improved anywhere in this note, and no ceiling at the scale
$\Tlin=C(t,U)\Omega$ is proved or disproved. The single exception is the
$\Ueff=0$ bond-SSH case, where Theorem~\ref{thm:canon} replaces the divergent
$t^2/\Ueff$ rate of \cite{sshbond3} by a finite bound. The obstructions are statements
about specified classes of arguments, as spelled out in
Remarks~\ref{rem:class}, \ref{rem:notclass} and \ref{rem:smooth}.
\item Theorems~\ref{thm:witness} and~\ref{thm:torus} concern a conditional coefficient,
not the Gibbs state. They do not show that the conditional pair correlations of the
witness field are large, and they do not show that such a field has appreciable weight
under any static or annealed measure.
\item Theorem~\ref{thm:amp} and Corollary~\ref{cor:fixture} concern a single-impurity
Green function along one history. They do not imply a divergent averaged propagator,
and they are not evidence that the impurity description is inappropriate for the
Holstein model.
\item Theorems~\ref{thm:gc}, \ref{thm:canon} and \ref{thm:routes} concern the pure
bond-SSH model with $\Ueff=0$ on square tori of side at least four. They do not
transfer to $\Ueff>0$, they say nothing at temperatures below $1024(|t|+\Gb)$, and
they are improvements only at $\Ueff=0$: the finite bound \eqref{eq:tpair} holds where
the rate $t^2/\Ueff$ of \cite{sshbond3} diverges, and at $t=0$ it lowers that paper's
threshold $2^{13}\Gb$ to $2^{10}\Gb$, which is a statement about the high-temperature
regime and not about the linear scale. They say nothing about the low-temperature
two-dimensional regime where pair correlations would be algebraic rather than
exponential.
\end{enumerate}

\section{Open problems}
\label{sec:open}

\begin{enumerate}[label=(\arabic*),leftmargin=2.2em]
\item \emph{The linear step.} Prove or refute
$k_B\Tpair\le C(t,U)\Omega$ at fixed $t,U$ and regular density as
$\Omega\downarrow0$. The obstructions here remove three specific mechanisms. What
remains open after them is a source-connected resummation that absorbs the connected
dynamic cumulant of Theorem~\ref{thm:witness} into a new reference, together with the
static estimate that controls how often near-Fermi resonant fields occur with
appreciable weight.
\item \emph{The annealed route.} Establish a volume-uniform bound on
$\E[\det(1+U_v)^2\mathcal{W}_v]/\E[\det(1+U_v)^2]$ for the relevant Wick minors or
polymer activities near non-separated clusters. Theorem~\ref{thm:amp} shows that the
numerator cannot be estimated history by history; adjugate minors remain bounded in
each finite dimension, and whether that structure is enough is open.
\item \emph{The Ward route.} Lemma~\ref{lem:fsum} and the identities of
Remark~\ref{rem:ward} are finite-matrix statements. To reach a ceiling one would need
the infinite-Matsubara limit of those identities, the identification of the exact
singlet kernel with the nonnegative expression of
Remark~\ref{rem:ward}, an averaged leakage estimate that does not replace $\alpha$ by
$\E\alpha$, and an operator-norm or weighted-resolvent version of a diagonal estimate.
None of these is available.
\item \emph{Dimer and Born-Oppenheimer extensions.} Extend the exact static landscape
of Remark~\ref{rem:dimer} to a volume-uniform quantum free-energy or coherence bound
with a controlled Born-Oppenheimer error expansion in $\Omega/t$. A finite dimer has no
pair-order transition, and a crossover defined from a single observable cannot be
relabelled $\Tpair$.
\item \emph{Reflection and infrared arguments.} The free optical covariance carries no
momentum-dependent stiffness, and ordinary spatial reflection positivity fails in
balanced canonical Holstein sectors; grand-canonical or particle-hole modified
reflections would need their own proof. A free Debye-Waller average cannot be factored
out of an interacting Lang-Firsov trace.
\item \emph{The $\Ueff=0$ model below the high-temperature regime.} Theorems
\ref{thm:gc} and~\ref{thm:canon} stop at $\beta|t|+\beta\Gb=2^{-10}$. Theorem~\ref{thm:routes}(iii) and the exact identities of
Theorem~\ref{thm:routes}(i) together show that the remaining regime is not a matter of
constants:
any estimate there must be joint in $\beta t$ and $\beta\Gb$, and it must recover
algebraic tails at large hopping.
\end{enumerate}

\section*{Acknowledgments and disclosures}

AI agents generated the proofs, code and drafts under D. Rens' direction, and
separate AI agents carried out the independent reviews.

This preprint has not been peer reviewed.

\section*{Data availability}

Every constant quoted in this note is reproduced by a replay suite that accompanies
this preprint: exact rational arithmetic for all rational comparisons and polynomial
identities, outward-rounded interval arithmetic for enclosures, and a separately
written independent implementation for each certificate, together with the
fixed-parameter fixtures of Corollary~\ref{cor:fixture} and the finite-matrix
diagnostics of the Remarks above. The finite-matrix and floating-point records
are numerical sanity checks and are labelled as such; they are not part of any proof.
The suite is deposited as the replay supplement of this preprint, \DOI{\SUPPDOI},
together with its source records at \url{\REPOURL}; this preprint is archived under
\DOI{\PAPERDOI}. No external data set is used anywhere in this work.

\bibliographystyle{plain}
\bibliography{refs}
\end{document}
